\colhead{学生练习}{第1—2题\quad 每题10分}
\exercise{1}{设全集 $U=\{-2,-1,0,1,2,3,4\}$，集合 $A=\{x\in\mathbb Z\mid-1\le x\le2\}$，$B=\{0,2,4\}$。\subq{1}{用列举法表示 $A$；}\subq{2}{求 $A\cap B$ 和 $A\cup B$；}\subq{3}{求 $\complement_U B$，并分别判断 $0\in A$、$\{0,2\}\subseteq A$ 是否成立。}}
\exercise{2}{设 $U=\mathbb R$，$A=\{x\mid x\le1\}$，$B=\{x\mid t<x<4\}$，其中 $t<4$。\subq{1}{当 $t=-2$ 时，求 $A\cap B$、$A\cup B$ 和 $\complement_U A$；}\subq{2}{若 $A\cup B=(-\infty,4)$，求 $t$ 的取值范围。}}
\newpage
\colhead{学生练习}{第3—4题\quad 每题10分}
\exercise{3}{设 $x\in\mathbb R$，分别判断下列两个条件之间的充分、必要关系，并说明理由。\subq{1}{$p$：$x=1$；$q$：$x^2=1$。判断 $p$ 是 $q$ 的什么条件；}\subq{2}{$r$：$x^2\le4$；$s$：$|x|\le2$。判断 $r$ 是 $s$ 的什么条件。}}
\exercise{4}{已知命题 $p$：$\forall x\in[0,+\infty)$，$x^2+1>0$；命题 $q$：$\exists x\in\mathbb R$，$x^2+1<0$。\subq{1}{写出命题p的否定；}\subq{2}{写出命题q的否定；}\subq{3}{判断p、q的真假，并说明理由。}}
\newpage
\colhead{学生练习}{第5—6题\quad 每题10分}
\exercise{5}{已知 $a>b>0$，$c<0$。\subq{1}{比较 $ac$ 与 $bc$、$\dfrac1a$ 与 $\dfrac1b$ 的大小；}\subq{2}{比较 $a^2c$ 与 $b^2c$ 的大小；}\subq{3}{“对任意 $d\in\mathbb R$，均有 $ad^2>bd^2$”是否正确？若不正确，给出反例，并写出对所有实数d都成立的大小关系及等号成立的条件。}}
\exercise{6}{解下列关于x的不等式，第（3）问中 $m\in\mathbb R$。\subq{1}{$-x^2+3x+4>0$；}\subq{2}{$(2x-1)^2\le9$；}\subq{3}{$(x-m)(x-1)\le0$。}}
\newpage
\colhead{学生练习}{第7—8题\quad 每题10分}
\exercise{7}{求下列最值，并写出取到最值时自变量的值。\subq{1}{当 $x>1$ 时，求 $x+\dfrac9{x-1}$ 的最小值；}\subq{2}{当 $0<x<8$ 时，求 $\sqrt{x(8-x)}$ 的最大值。}}
\exercise{8}{已知 $m\in\mathbb R$。\subq{1}{若对任意 $x\in\mathbb R$，都有 $x^2-4x+m\ge0$，求m的取值范围；}\subq{2}{若将第（1）问中的“$\ge0$”改为“$>0$”，求m的取值范围。}}
\newpage
\colhead{学生练习}{第9—10题\quad 每题10分}
\exercise{9}{回答下列问题，每问都说明理由。\subq{1}{$y=\sqrt{x^2}$ 与 $y=|x|$ 是否表示同一个函数？}\subq{2}{$y=x$ 与 $y=\dfrac{x^2}{x}$ 是否表示同一个函数？}\subq{3}{设 $A=\{-1,0,1\}$，$B=\{0,1,2\}$。对应 $x\mapsto x^2$ 能否构成从A到B的函数？它的值域是否等于B？}}
\exercise{10}{求下列函数的定义域。\subq{1}{$f(x)=\dfrac{\sqrt{3-x}}{x-1}$；}\subq{2}{$g(x)=\sqrt{x-1}+\dfrac1{x-3}$。}}
\newpage
\colhead{学生练习}{第11—12题\quad 每题10分}
\exercise{11}{已知 $f(x)=\begin{cases}x+2,&x\le0,\\x^2,&x>0.\end{cases}$\subq{1}{求 $f(0)$、$f(-1)$ 和 $f(f(-1))$；}\subq{2}{函数g的对应关系如下表，求 $g(g(2))$。\par\smallskip\begin{tabular}{c|ccc}$x$&1&2&3\\\hline$g(x)$&3&1&2\end{tabular}}\subq{3}{若对任意 $x\in\mathbb R$ 都有 $h(x+2)=2x-1$，求 $h(3)$，并求 $h(x)$ 的解析式。}}
\exercise{12}{函数 $y=f(x)$ 的图象为折线ABC，$A(0,2)$、$B(1,0)$、$C(3,2)$。\readinggraph\subq{1}{写出f的定义域和值域；}\subq{2}{求 $f(0)$ 和 $f(f(0))$；}\subq{3}{求不等式 $f(x)\le1$ 的解集。}}
\newpage
\colhead{学生练习}{第13—14题\quad 每题10分}
\exercise{13}{已知函数 $f(x)=x^2-2x$，$x\in[-1,3]$。\subq{1}{写出f的单调递减区间和单调递增区间；}\subq{2}{求f的最大值、最小值，并写出取得各最值时x的所有取值；}\subq{3}{写出f的值域。}}
\exercise{14}{完成下列问题。\subq{1}{用定义证明函数 $f(x)=x^2$ 在 $[0,+\infty)$ 上单调递增；}\subq{2}{若函数 $h(x)=x^2-2ax$ 在 $(-\infty,3]$ 上单调递减，求实数a的取值范围。}}
\newpage
\colhead{学生练习}{第15—16题\quad 每题10分}
\exercise{15}{已知f是定义在 $\mathbb R$ 上的奇函数，当 $x\ge0$ 时，$f(x)=x^2+2x$。\subq{1}{求 $f(0)$ 和 $f(-2)$；}\subq{2}{求f在 $(-\infty,0)$ 上的解析式；}\subq{3}{将 $f(-3)$、$f(-1)$、$f(2)$ 按从小到大排列。}}
\exercise{16}{已知g的定义域为 $[-2,2]$，设 $f(x)=g(x)+g(-x)$。当 $x\in[0,2]$ 时，$f(x)=2x-x^2$。\subq{1}{判断f的奇偶性，并证明；}\subq{2}{求f在 $[-2,0)$ 上的解析式；}\subq{3}{求f在 $[-2,2]$ 上的最大值、最小值，并写出所有最值点的横坐标。}}
\newpage
\colhead{学生练习}{第17—18题\quad 每题10分}
\exercise{17}{计算下列各式，写出必要步骤。\subq{1}{$32^{\frac25}+27^{-\frac23}$；}\subq{2}{$16^{\frac34}+8^{-\frac23}$；}\subq{3}{$\sqrt{(-3)^2}+\sqrt[3]{-8}$。}}
\exercise{18}{已知幂函数 $f(x)=x^\alpha$ 的图象经过点 $\left(4,\dfrac1{16}\right)$。\subq{1}{求f的解析式及定义域；}\subq{2}{判断f的奇偶性，并写出它的单调区间及增减性；}\subq{3}{用定义证明f在 $(0,+\infty)$ 上单调递减；}\subq{4}{分别写出函数 $h(x)=x^{3/2}$、$k(x)=x^{-1/2}$ 的实数定义域。}}
\newpage
\colhead{学生练习}{第19—20题\quad 每题10分}
\exercise{19}{已知指数函数 $f(x)=a^x$ 的图象经过点 $(2,4)$，其中 $a>0$ 且 $a\ne1$，设 $g(x)=a^{|x-1|}$。\subq{1}{求a的值；}\subq{2}{写出g的单调区间及相应的增减性，并求g的值域；}\subq{3}{解不等式 $g(x)\le8$。}}
\exercise{20}{某种物质的剩余量 $M(t)$（g）满足 $M(t)=100a^t$，其中t为经过的时间（年），$t\ge0$，$a>0$ 且 $a\ne1$。已知两年后剩余25 g。\subq{1}{求a和 $M(t)$ 的解析式；}\subq{2}{判断 $M(t)$ 的单调性，写出值域，并说明剩余量能否等于0；}\subq{3}{当剩余量不少于12.5 g时，求t的取值范围。}}
\newpage
